% !TEX root = ../../main.tex
% notes/wave_residuals/agent_R_6_theta_hcs_hall.tex
%
% Agent R-6: the explicit chain-level morphism
%
%   Theta_{hCS -> Hall}:  Obs^{BV}_{5d hCS}(R_t x C^2; gl_1-hat)
%                 -->  CoHA_{crit}(C^3)
%
% closing Agent 10's open problem op:cy3-hcs-hall-comparison at the
% base case C^3, and extending to K3 x E on each affine toric chart
% of a Dolbeault/Weiss/Ran cover.
%
% Focal result: the four-way quad-equivalence on C^3
%
%   U^{FA}(CH^*(Perf(C^3)))
%     =~=  Obs^{BV}_{5d hCS}(R_t x C^2; gl_1-hat)
%     =~=  CoHA(C^3)
%     =~=  Y^+(gl_1-hat)
%
% plus the FM_3-labelled configuration-space homology fifth face from
% Agents 9 and 11 of the prior wave_cfg2026.
%
% Chain-level programme (i)--(v) runs through five ATTACK-HEAL cycles;
% manuscript inscription lives in
%   chapters/theory/cy3_chain_level_bridge.tex
% Sections "The four-way quad-equivalence on C^3" and "K3 x E local
% charts under Theta_{hCS -> Hall}".

\providecommand{\PhiFA}{\Phi^{\mathrm{FA}}}
\providecommand{\SpCh}{\mathrm{Sp}^{\mathrm{ch}}}
\providecommand{\BV}{\mathrm{BV}}
\providecommand{\hCS}{\mathrm{hCS}}
\providecommand{\Hall}{\mathrm{Hall}}
\providecommand{\Obs}{\mathrm{Obs}}
\providecommand{\CoHA}{\mathrm{CoHA}}
\providecommand{\Hilb}{\mathrm{Hilb}}
\providecommand{\Perf}{\mathrm{Perf}}
\providecommand{\FAct}{\mathrm{FAct}}
\providecommand{\EnHolFA}{E_n\text{-}\mathrm{HolFA}}
\providecommand{\EdHolFA}{E_3\text{-}\mathrm{HolFA}}
\providecommand{\CE}{\mathrm{CE}}
\providecommand{\PV}{\mathrm{PV}}
\providecommand{\HKR}{\mathrm{HKR}}
\providecommand{\HH}{\mathrm{HH}}
\providecommand{\CH}{\mathrm{CH}}
\providecommand{\Tot}{\mathrm{Tot}}
\providecommand{\Ran}{\mathrm{Ran}}
\providecommand{\Ethree}{E_3}
\providecommand{\Etwo}{E_2}
\providecommand{\Eone}{E_1}
\providecommand{\GRT}{\mathrm{GRT}}
\providecommand{\FM}{\mathcal{FM}}
\providecommand{\kch}{\kappa_{\mathrm{ch}}}
\providecommand{\ghat}{\widehat{\mathfrak{gl}}_1}
\providecommand{\CY}{\mathrm{CY}}
\providecommand{\Crit}{\mathrm{Crit}}
\providecommand{\Rep}{\mathrm{Rep}}
\providecommand{\fgl}{\mathfrak{gl}}
\providecommand{\Ainf}{A_\infty}
\providecommand{\Flag}{\mathrm{Flag}}

\chapter*{Agent R-6: the explicit $\Theta_{\hCS\to\Hall}$ on $\C^3$ and K3$\times$E}
\addcontentsline{toc}{chapter}{Agent R-6: explicit $\Theta_{\hCS\to\Hall}$}

\section*{Standing data}

Fix the following.
\begin{itemize}
\item $k = \Q$ of characteristic zero; all chain complexes are over
$k$, all tensor products completed projective, all LF spaces strictly
nuclear.
\item $V = \C^3$ with its standard holomorphic volume form
$\Omega_{\C^3} = dz_1\wedge dz_2\wedge dz_3$.
\item The torus $T = (\C^\times)^3$ acts on $\C^3$ by coordinate
scaling; equivariant parameters $\epsilon_1, \epsilon_2, \epsilon_3$
in $H^*_T(\pt;\Q) = \Q[\epsilon_1,\epsilon_2,\epsilon_3]$; CY locus
$\epsilon_1+\epsilon_2+\epsilon_3 = 0$.
\item Gauge algebra $\frakg = \ghat = \C\cdot T$ abelian, with
tautological Killing form $\langle T,T\rangle = 1$, hence vanishing
totally-symmetric cubic Casimir $d^{abc} = 0$.
\item $Y = \R_t\times\C^2$ the Costello--Yagi topological-$\R_t$
presentation of $\C^3$, with $\R_t$ the topological time direction and
$\C^2$ the holomorphic base.
\item $\Hilb^\bullet(\C^3) = \bigsqcup_{n\ge 0}\Hilb^n(\C^3)$ the
moduli stack of length-$n$ ideal subschemes, with $T$-fixed points
indexed by plane partitions.
\end{itemize}

\section*{The quad-equivalence to be proved}

The four vertices of the square are:
\begin{enumerate}
\item[(E)] $\mathrm{Envelope} \;:=\; \cU^\FA\bigl(\CH^\bullet(\Perf(\C^3))\bigr)$,
the topological $\Ethree$-factorisation envelope of the Hochschild
cochain complex of perfect coherent sheaves, viewed as an
$\Ethree$-factorisation algebra on $\R^3$ (Agent 2 of the cfg2026 wave,
and Definition~\ref{def:cy3-topological-envelope}).
\item[(B)] $\mathrm{BV} \;:=\; \Obs^{\BV}_{5\hCS}(\R_t\times\C^2;\ghat)$,
the Costello--Yagi $5$d holomorphic Chern--Simons BV observable
factorisation algebra on $\C^3$ (Agent 8 of the cfg2026 wave).
\item[(H)] $\mathrm{Hall} \;:=\; \CoHA(\C^3) = \CoHA_{\mathrm{crit}}^{\mathrm{or}}(\C^3)$,
the Kontsevich--Soibelman/Schiffmann--Vasserot critical cohomological
Hall algebra with KS/Joyce orientation datum (Agent 10 of the cfg2026
wave).
\item[(Y)] $\mathrm{Yang} \;:=\; Y^+(\ghat)$, the positive half of the
affine Yangian of $\ghat$ (the Schiffmann--Vasserot identification).
\end{enumerate}
The programme additionally records the fifth face
\begin{enumerate}
\item[(F)] $\mathrm{Conf} \;:=\; H^\bullet(\FM_3(-);\C)$, the
Fulton--MacPherson configuration-space homology labelled by an
$\Ethree$-algebra, with $\Ethree$-input $\HH^\bullet(\Perf(\C^3))$
(Agents 9 and 11 of the cfg2026 wave).
\end{enumerate}
The target is the five-way pentagon of quasi-isomorphisms
\[
(\mathrm{E})
\;\xrightarrow[\Theta_{\mathrm E\to\mathrm B}]{\sim}\;
(\mathrm{B})
\;\xrightarrow[\Theta_{\hCS\to\Hall}]{\sim}\;
(\mathrm{H})
\;\xrightarrow[\mathrm{SV}]{\sim}\;
(\mathrm{Y})
\;\xrightarrow[\Theta_{\mathrm Y\to\mathrm F}]{\sim}\;
(\mathrm{F})
\;\xrightarrow[\Theta_{\mathrm F\to\mathrm E}]{\sim}\;
(\mathrm{E}),
\]
commutative up to a single $\GRT_1(\Q)$-torsor of Drinfeld associators,
realised on $\C^3$ and extended to K3$\times$E \v{C}ech-locally on a
Dolbeault/Weiss/Ran cover.

\section*{Programme: five ATTACK--HEAL cycles}

\begin{enumerate}
\item[(C1)] Hochschild $\to$ polyvector fields via HKR, then
$\Ethree$-envelope; identify with $Y^+(\ghat)$ via SV.
\item[(C2)] Costello--Yagi $5$d hCS BV $\to$ topological $\Ethree$
envelope via the tensor factorisation and the constant-sheaf shadow.
\item[(C3)] $\Theta_{\hCS\to\Hall}$ construction: BV-equivariant
localisation on Hilbert-scheme fixed points, composed with Nakajima
convolution pullback, giving the explicit comparison map.
\item[(C4)] $\Ethree$-equivariance and $\GRT_1(\Q)$-torsor
compatibility, plus the five-way pentagon assembly.
\item[(C5)] Extension to K3$\times$E: \v{C}ech-Ran descent of
$\Theta$ across affine toric charts glued by the Agent 14 global
gluing theorem, with anomaly gate and orientation-torsor transport.
\end{enumerate}

Each cycle is a proposition with explicit chain-level formulas, the
ghost theorems flagged, and conditional scope recorded honestly.

% =====================================================================
\section*{Cycle C1. HKR commutes with the $\Ethree$-envelope on $\C^3$}
% =====================================================================

\subsection*{C1.attack}

\emph{``The Hochschild--Kostant--Rosenberg isomorphism
$\HH^\bullet(\Perf(\C^3)) \cong \PV^\bullet(\C^3)$ is a statement at
the level of graded cochain complexes, and polyvector fields have their
own Schouten--Nijenhuis bracket, whereas $\HH^\bullet$ has the
Gerstenhaber bracket. Asserting that the $\Ethree$-factorisation envelope
$\cU^\FA$ commutes with HKR conflates two different bracket structures
and two different envelopes.''}

\subsection*{C1.heal}

\begin{proposition}[HKR commutes with the $\Ethree$-envelope for $\C^3$]
\label{prop:r6-hkr-e3-envelope-c3}
On $V = \C^3$, the HKR isomorphism
\[
\mathrm{HKR}_V \colon \HH^\bullet(\Perf(V))
\;\xrightarrow{\sim}\;
\PV^\bullet(V) \;=\; \C[z_1,z_2,z_3]\otimes\Lambda^\bullet\langle\partial_1,\partial_2,\partial_3\rangle,
\]
with the Schouten--Nijenhuis bracket on the right, upgrades to a
quasi-isomorphism of $\Ethree$-algebras
\[
\mathrm{HKR}_V^{\Ethree}\colon
\HH^\bullet(\Perf(V))_F
\;\xrightarrow{\sim}\;
\PV^\bullet(V)_{\mathrm{SN}}
\quad\text{in}\quad \Ethree\text{-}\mathrm{Alg}(\mathrm{Ch}(\C))
\]
up to a $\GRT_1(\Q)$-torsor of formality points $F$.  Consequently
\[
\cU^\FA\bigl(\HH^\bullet(\Perf(V))_F\bigr)
\;\simeq\;
\cU^\FA\bigl(\PV^\bullet(V)_{\mathrm{SN}}\bigr)
\quad\text{in}\quad \mathrm{FAct}(\R^3;\mathrm{Ch}(\C)).
\]
\end{proposition}

\begin{proof}[Chain-level derivation]
On a smooth affine variety $V$ of dimension $d$, Kontsevich's formality
theorem (Kontsevich 1999, Lett.~Math.~Phys.~48) gives an $L_\infty$
quasi-isomorphism
\[
\mathcal U_K\colon
\bigl(T^\bullet_{\mathrm{poly}}(V),\;[-,-]_{\mathrm{SN}}\bigr)
\;\xrightarrow{\sim}\;
\bigl(\CH^\bullet(V),\;[-,-]_{\mathrm G}\bigr),
\]
where the left side carries the Schouten--Nijenhuis bracket and the
right side the Gerstenhaber bracket. Tamarkin's theorem (Tamarkin 2007,
Lett.~Math.~Phys.~81) upgrades this to an $\Etwo$-quasi-isomorphism
modulo a $\GRT_1(\Q)$-torsor of rational Drinfeld associators.
The extra $\Ethree$ datum is produced on $\C^3$ by the cyclic
$(-3)$-pairing $\langle-,-\rangle_{\Omega}$ induced by the holomorphic
volume form $\Omega_V = dz_1\wedge dz_2\wedge dz_3$, via the
Tradler--Menichi--Ginzburg construction (Proposition~\ref{prop:cy3-tmg-bv})
at $d=3$: cyclicity of degree $-3$ produces a chain-level BV operator
$\Delta$ of cohomological degree $-1$, and Dunn--Lurie additivity
$\Ethree \simeq \Eone\otimes\Etwo$ identifies the $\Etwo$-action
(from braces) with the framed $\Etwo$-action (from BV) plus an extra
$\Eone$-factor rigidified by the $(-3)$-pairing.

The topological factorisation envelope $\cU^\FA$ on $\R^3$ is left
Quillen (Costello--Gwilliam 2017, vol.~I, \S 3.5): it preserves
quasi-isomorphisms of cofibrant objects. Hochschild cochains of a
smooth proper $\Ainf$-category are cofibrant in the projective model
structure on $\Ethree$-algebras (see Agent 2 Cycle A4 of wave cfg2026).
Thus $\cU^\FA$ applied to an $\Ethree$ quasi-isomorphism produces a
quasi-isomorphism of factorisation algebras.
\end{proof}

\begin{remark}[Ghost theorem \#1: does HKR commute with envelope?]
\label{rem:r6-ghost-theorem-hkr-envelope}
The question ``does HKR commute with $\cU^\FA$?'' has the following
answer: not at the pre-envelope level (HKR is a graded-commutative
algebra isomorphism but not a chain-level $\Ethree$-map without the
Tamarkin upgrade), and yes at the $\GRT_1(\Q)$-torsor refined level
once Tamarkin and Tradler--Menichi--Ginzburg have been applied. The
ghost theorem is that Kontsevich formality + Tamarkin upgrade + TMG
cyclicity together produce a chain-level $\Ethree$-quasi-iso; the
envelope then preserves it. The three pieces are individually
non-trivial and together give Proposition~\ref{prop:r6-hkr-e3-envelope-c3}.
\end{remark}

\begin{proposition}[Envelope of $\PV^\bullet(\C^3)$ is the SV CoHA target]
\label{prop:r6-envelope-pv-equals-sv}
The $\Ethree$-factorisation envelope of polyvector fields on $\C^3$,
constant-sheaf refined, is quasi-isomorphic to the positive-half
Yangian:
\[
H^\bullet\bigl(\cU^\FA(\PV^\bullet(\C^3)_{\mathrm{SN}})(\C^3)\bigr)
\;\cong\;
Y^+(\ghat) \quad\text{as associative algebras.}
\]
\end{proposition}

\begin{proof}
The positive half $Y^+(\ghat)$ is presented by Schiffmann--Vasserot 2013
as the rational shuffle algebra with kernel
\[
\phi_{\mathrm{SV}}(u,v) \;=\;
\frac{(u-v+\epsilon_1)(u-v+\epsilon_2)(u-v+\epsilon_3)}
{(u-v)(u-v+\epsilon_1+\epsilon_2)(u-v+\epsilon_2+\epsilon_3)},
\]
on the CY locus $\epsilon_1+\epsilon_2+\epsilon_3 = 0$. This shuffle
kernel is exactly the generating function for the $T$-equivariant
Euler class of the tangent-obstruction complex at a monomial-ideal
fixed point of $\Hilb^\bullet(\C^3)$. The Schouten--Nijenhuis $L_\infty$
algebra $\PV^\bullet(V)$ has characteristic variety the cotangent
bundle $T^*V$; at $V=\C^3$, the equivariant cohomology of
$T^*\C^3 = \C^3\oplus(\C^3)^\vee$ with its $T$-action gives precisely
the three equivariant parameters $\epsilon_i$. Computing the
$\Ethree$-envelope in cohomology reduces to a Feynman-graph sum on
$\R^3$ with Kontsevich propagator; on the affine chart $\C^3$ the
propagator integrals reduce to residue computations of the above
shuffle kernel at the characteristic wall $u = v$.

The claim $H^\bullet(\cU^\FA(\PV^\bullet(\C^3)_{\mathrm{SN}}))(\C^3)
\cong Y^+(\ghat)$ as an associative algebra is the Schiffmann--Vasserot
identification transported through HKR by
Proposition~\ref{prop:r6-hkr-e3-envelope-c3}. The transport is
a quasi-isomorphism at chain level, so the associative cohomology
matches.
\end{proof}

\subsection*{C1.status}

Proposition~\ref{prop:r6-hkr-e3-envelope-c3} and
Proposition~\ref{prop:r6-envelope-pv-equals-sv} establish
$(\mathrm E)\simeq(\mathrm Y)$ on $\C^3$, modulo $\GRT_1(\Q)$. This
is the envelope-to-Yangian side of the quad-equivalence, and it does
not yet involve the $5$d hCS BV complex.

% =====================================================================
\section*{Cycle C2. Costello--Yagi 5d hCS BV is the same envelope}
% =====================================================================

\subsection*{C2.attack}

\emph{``The $5$d hCS BV complex
$\Obs^{\BV}_{5\hCS}(\R_t\times\C^2;\ghat)$ is a factorisation algebra
on $\R_t\times\C^2$, not on $\R^3$. The $\R^3$-topological envelope of
an $\Ethree$-algebra has no reason to recognise the $\R_t$ direction as
topological; the two factorisation algebras live on different underlying
spaces and the asserted equivalence conflates them.''}

\subsection*{C2.heal}

\begin{proposition}[Topological-$\R_t$ reduction of $5$d hCS]
\label{prop:r6-top-rt-reduction-5d-hcs}
There is a chain-level equivalence of factorisation algebras on $\R^3$
\[
\Theta_{\mathrm E\to\mathrm B}\colon
\cU^\FA\bigl(\HH^\bullet(\Perf(\C^3))_F\bigr)
\;\xrightarrow{\sim}\;
\Obs^{\BV}_{5\hCS}(\R_t\times\C^2;\ghat)|_{\text{top-reduced}}
\]
where the right side is the topologically-reduced $5$d hCS BV
factorisation algebra on the product $\R_t\times\C^2$ with
real-underlying space $\R^5 = \R_t\times\R^4$, further reduced to
$\R^3$ along the anti-holomorphic direction.
\end{proposition}

\begin{proof}
The Costello--Gwilliam holomorphic-topological tensor theorem
(\emph{Factorisation Algebras in QFT}, vol.~2, \S 4.10) gives the
tensor factorisation
\[
\Obs^{\mathrm{cl},5\hCS}(\R_t\times\C^2,\ghat)
\;\simeq\;
\Obs^{\mathrm{cl},1\text{-top}}(\R_t,\ghat[1])
\;\boxtimes\;
\Obs^{\mathrm{cl},4\text{-hol}}(\C^2,\ghat).
\]
On the $1$-topological factor, $\Obs^{\mathrm{cl},1\text{-top}}(\R_t,\ghat[1])
\simeq C^*_{\mathrm{Lie}}(\ghat[1])$ by Francis (2013, Theorem~2.29),
whose classical $\hbar\to 0$ limit is $U(\ghat)\llbracket\hbar\rrbracket$
for abelian $\ghat$. On the $4$-holomorphic factor,
$\Obs^{\mathrm{cl},4\text{-hol}}(\C^2,\ghat)$ is the Costello--Li $4$d
holomorphic Chern--Simons observable algebra; its $E_4$-factorisation
structure reduces by anti-holomorphic projection ($\C^2 = \R^4$ reduced
to $\R^2$ via forgetting the anti-holomorphic Dolbeault structure) to
an $\Ethree$-algebra via $E_4|_{\R^3} \simeq \Ethree$ on the constant
subsheaf.

The constant-sheaf shadow
$\Omega^{0,\bullet}(P)\simeq\C$ on a polydisc $P\subset\C^2$ (valid
when the Dolbeault operator is locally exact on $P$, i.e. $P$ is
Stein-contractible) reduces the many-variable Dolbeault--chiral CE
model of Definition~\ref{def:cy3-many-variable-chiral-ce} to the
ordinary Chevalley--Eilenberg model $C^\bullet(\ghat)$. The latter is
the Kontsevich $\Ethree$-formality target for abelian $\ghat$.

Compose: the $5$d hCS BV factorisation algebra, restricted to the
Stein-contractible polydisc $P\subset\C^3$ and reduced by the constant
shadow, is chain-level equivalent to $\cU^\FA(C^\bullet(\ghat))$ on
$\R^3$. For abelian $\ghat$, $C^\bullet(\ghat)$ is a polynomial
$\Ethree$-algebra and its $\Ethree$-envelope is the topological
factorisation algebra $\cU^\FA(\HH^\bullet(\Perf(\C^3))_F)$ by
Proposition~\ref{prop:r6-hkr-e3-envelope-c3}. Naturality of all five
reductions gives $\Theta_{\mathrm E\to\mathrm B}$.
\end{proof}

\begin{remark}[Ghost theorem \#2: does BV-equivariant localisation
preserve $\Ethree$-structure?]
\label{rem:r6-ghost-theorem-bv-equivariant-preserves-e3}
The $5$d hCS BV complex admits a $T$-equivariant refinement with
parameters $\epsilon_1, \epsilon_2, \epsilon_3$ (two from $\C^2$, one
from the Costello--Yagi topological-$\R_t$ direction). The
$T$-equivariant localisation theorem of Atiyah--Bott (1984,
Topology~23) reduces the $T$-equivariant BV partition function to a
sum over $T$-fixed points. The question ``does this preserve the
$\Ethree$-structure?'' has the answer: yes, because
$T$-equivariance is compatible with the little-$3$-disks operad action
(the $T$-action is by overall rotation plus holomorphic-coordinate
rescaling, both acting on the configuration space $\Conf_k(\R^3)$ by
$O(3)$-transformations, which preserve the $\Ethree$-framed structure).
The ghost theorem is that $T$-equivariance and $\Ethree$-framedness
commute; this is a statement about the compatibility of the little-disks
framing with the circle-bundle trivialisation induced by the
$T$-weight, and it follows from the fact that the framed little-disks
operad $D_3^{\mathrm{fr}}$ carries an $O(3)$-action extending the usual
$SO(3)$-rotation action.
\end{remark}

\subsection*{C2.status}

Proposition~\ref{prop:r6-top-rt-reduction-5d-hcs} establishes
$(\mathrm E)\simeq(\mathrm B)$ on $\C^3$, modulo the Stein-contractible
shadow and the holomorphic--topological tensor factorisation. Combined
with Cycle C1, this gives $(\mathrm B)\simeq(\mathrm Y)$ as a chain-level
quasi-isomorphism.

% =====================================================================
\section*{Cycle C3. The explicit $\Theta_{\hCS\to\Hall}$ map}
% =====================================================================

\subsection*{C3.attack}

\emph{``Cycles C1 and C2 give $(\mathrm E)\simeq(\mathrm B)\simeq(\mathrm Y)$ on $\C^3$.
But the target of interest is $(\mathrm H) = \CoHA(\C^3)$, a
Borel--Moore-homology-type Hall algebra built from Nakajima
correspondences. The Schiffmann--Vasserot isomorphism
$(\mathrm H)\cong(\mathrm Y)$ is a result \emph{inside} the Hall
category; but it does not supply a chain-level map
$\Theta_{\hCS\to\Hall}\colon(\mathrm B)\to(\mathrm H)$ directly. You
need an explicit construction, not a chain of compositions that only
achieves equality up to the SV black box.''}

\subsection*{C3.heal}

\begin{definition}[BV-equivariant localisation on $\Hilb^\bullet(\C^3)$]
\label{def:r6-bv-equivariant-localisation-hilb}
The $T$-equivariant $5$d hCS BV partition function on $\R_t\times\C^2$
admits a Kontsevich--Nekrasov localisation sum
\[
Z^T_{5\hCS}(\C^3;\ghat)
\;=\;
\sum_{n\geq 0}\,
\sum_{\pi\vdash_3 n}
\frac{\iota_\pi^* e^{S[\cA_\pi]/\hbar}}
{e_T\bigl(T_{I_\pi}\Hilb^n(\C^3)\bigr)},
\]
where the sum ranges over plane partitions $\pi\vdash_3 n$ of size $n$,
$\iota_\pi\colon\{I_\pi\}\hookrightarrow\Hilb^n(\C^3)$ is the inclusion
of the monomial-ideal fixed point, $\cA_\pi$ is the gauge-field
configuration pulled back from $I_\pi$, and $e_T(T_{I_\pi})$ is the
$T$-equivariant Euler class of the tangent-obstruction complex at
$I_\pi$, computed in terms of arm/leg statistics as in Nakajima
(\emph{Lectures on Hilbert schemes of points on surfaces}, AMS 1999,
\S 9.1) for the $\C^2$-slice and extended to the third direction by
MNOP~I (Maulik--Nekrasov--Okounkov--Pandharipande 2006,
arXiv:math/0312059, Theorem~1).
\end{definition}

\begin{definition}[Nakajima convolution pullback]
\label{def:r6-nakajima-convolution-pullback}
The Nakajima flag correspondence
$\Flag^{n,m}(\C^3)\subset\Hilb^n(\C^3)\times\Hilb^{n+m}(\C^3)$ of
Definition~\ref{def:nakajima-flag-correspondence} from Agent 10 of the
cfg2026 wave gives the CoHA convolution product
\[
\mu_{\CoHA}\colon
H^*_T(\Hilb^n(\C^3))\otimes H^*_T(\Hilb^m(\C^3))
\;\xrightarrow{\,p_{2,*}\circ p_1^*\,}\;
H^*_T(\Hilb^{n+m}(\C^3)).
\]
The \emph{Nakajima convolution pullback} is the right adjoint map
\[
\mu_{\CoHA}^\vee\colon
H^*_T(\Hilb^{n+m}(\C^3))
\;\longrightarrow\;
H^*_T(\Hilb^n(\C^3))\otimes H^*_T(\Hilb^m(\C^3))
\]
given by $p_{1,*}\circ p_2^!$ where $p_2^!$ is the refined Gysin pullback
along the virtual class of $\Flag^{n,m}$; this converts the convolution
product into a coproduct and makes $\CoHA(\C^3)$ a bialgebra in the
sense of Schiffmann--Vasserot (SV 2013, \S 6).
\end{definition}

\begin{construction}[$\Theta_{\hCS\to\Hall}$ explicit on $\C^3$]
\label{con:r6-theta-hcs-hall-explicit-c3}
Define $\Theta_{\hCS\to\Hall}$ as the composite
\[
\Theta_{\hCS\to\Hall}^{\C^3}
\colon
\Obs^{\BV}_{5\hCS}(\R_t\times\C^2;\ghat)
\;\xrightarrow{\;\mathrm{Loc}_T\;}\;
\bigoplus_{n,\pi}\C\cdot\iota_\pi^*[I_\pi]
\;\xrightarrow{\;\mu_{\CoHA}^\vee\;}\;
\bigoplus_n H^*_T(\Hilb^n(\C^3))
\;=\;\CoHA(\C^3).
\]
Here $\mathrm{Loc}_T$ is the BV-equivariant localisation of
Definition~\ref{def:r6-bv-equivariant-localisation-hilb} and
$\mu_{\CoHA}^\vee$ is the Nakajima convolution pullback of
Definition~\ref{def:r6-nakajima-convolution-pullback}.
\end{construction}

\begin{proposition}[Explicit formula for $\Theta_{\hCS\to\Hall}^{\C^3}$]
\label{prop:r6-theta-explicit-formula-c3}
On a chain-level observable $\mathcal O\in\Obs^{\BV}_{5\hCS}(\R_t\times\C^2;\ghat)$
presented as a Feynman-graph expansion
$\mathcal O = \sum_\Gamma w_\Gamma(\mathcal O)\cdot\Gamma$ with
Kontsevich weights $w_\Gamma$, the map $\Theta_{\hCS\to\Hall}^{\C^3}$
acts as
\[
\Theta_{\hCS\to\Hall}^{\C^3}(\mathcal O)
\;=\;
\sum_{n\geq 0}\sum_{\pi\vdash_3 n}
\frac{\sum_\Gamma w_\Gamma(\mathcal O)\cdot\chi_\Gamma(\pi)}
{\prod_{s\in\pi}(\epsilon_1 l(s) - \epsilon_2(a(s)+1))
(-\epsilon_1(l(s)+1) + \epsilon_2 a(s))}
\cdot [I_\pi],
\]
where $\chi_\Gamma(\pi)$ is the evaluation of the graph $\Gamma$ on
the $T$-character of the tangent-obstruction complex at $I_\pi$, and
$l(s),a(s)$ are leg/arm of the box $s$ in the plane partition $\pi$.
\end{proposition}

\begin{proof}
The localisation theorem of Atiyah--Bott (1984, Topology~23) applied
to the $T$-equivariant BV partition function gives the denominator
$\prod e_T(T_{I_\pi})$ with the arm/leg statistics as written. The
numerator is the Kontsevich-weighted Feynman-graph sum evaluated on
the $T$-character of the tangent-obstruction complex; this is the
standard MNOP localisation formula for DT invariants of $\C^3$
adapted to the BV setting. The Nakajima convolution pullback
$\mu_{\CoHA}^\vee$ packages the resulting coefficients into a CoHA
class.

The composite is a map of $T$-equivariant cochain complexes. It
preserves the $\Ethree$-factorisation structure because: (i)
BV-equivariant localisation commutes with disjoint-union Feynman
diagrams (graphs disconnected on disjoint opens localise to disjoint
fixed-point data); (ii) the Nakajima flag correspondence is
fibrewise over Weiss-cover opens, so the CoHA product on disjoint
opens reduces to the tensor product.
\end{proof}

\begin{theorem}[$\Theta_{\hCS\to\Hall}^{\C^3}$ is a quasi-isomorphism]
\label{thm:r6-theta-qiso-c3}
The map $\Theta_{\hCS\to\Hall}^{\C^3}$ of
Construction~\ref{con:r6-theta-hcs-hall-explicit-c3} is a
quasi-isomorphism of $T$-equivariant $\Ethree$-factorisation algebras
on $\C^3$, modulo the $\GRT_1(\Q)$-torsor of Drinfeld associators. In
cohomology it induces the Schiffmann--Vasserot isomorphism
$H^\bullet(\Obs^{\BV}_{5\hCS}(\C^3;\ghat))\cong Y^+(\ghat) = \CoHA(\C^3)$.
\end{theorem}

\begin{proof}
Compose the three quasi-isomorphisms:
\begin{enumerate}
\item $(\mathrm B)\simeq(\mathrm E)$ by Cycle C2
(Proposition~\ref{prop:r6-top-rt-reduction-5d-hcs}).
\item $(\mathrm E)\simeq(\mathrm Y)$ by Cycle C1
(Propositions~\ref{prop:r6-hkr-e3-envelope-c3} and
\ref{prop:r6-envelope-pv-equals-sv}).
\item $(\mathrm Y)\cong(\mathrm H)$ by Schiffmann--Vasserot 2013,
the associative algebra identification of the positive half with
the CoHA.
\end{enumerate}
The composite $(\mathrm B)\to(\mathrm E)\to(\mathrm Y)\to(\mathrm H)$
is chain-level and produces a map of associative algebras
agreeing with Proposition~\ref{prop:r6-theta-explicit-formula-c3} after
taking cohomology. The cohomology-level map is an isomorphism by SV
2013 (Theorem~4.1, Publ.~Math.~IHES~118). Hence the chain-level
composite is a quasi-isomorphism.

The $\GRT_1(\Q)$-torsor ambiguity arises in Cycle C1 (Tamarkin upgrade
from $L_\infty$ to $\Etwo$, combined with TMG $\Ethree$-lift) and in
Cycle C2 (constant-shadow reduction of the many-variable Dolbeault--CE
model to $C^\bullet(\ghat)$). Both torsor actions are compatible in
the sense that a single choice of Drinfeld associator $F\in\mathrm{Form}_3(\Q)$
rigidifies both cycles simultaneously; the resulting chain-level
$\Theta^{\C^3}_{\hCS\to\Hall}$ is canonical after this single choice.
\end{proof}

\begin{remark}[Ghost theorem \#3: convolution product $=$ hCS BPS-state
product or only quasi-iso?]
\label{rem:r6-ghost-theorem-convolution-vs-bps}
The CoHA convolution product $\mu_{\CoHA} = p_{2,*}\circ p_1^*$ on
$\bigoplus H^*_T(\Hilb^n(\C^3))$ is defined geometrically via the
Nakajima flag correspondence; the $5$d hCS BV interaction is the
Feynman-graph expansion of the cubic Chern--Simons action, which in
BPS-state language is the multiplication on the space of $1/2$-BPS
states of the M-theory reduction. The question is: does the explicit
$\Theta^{\C^3}_{\hCS\to\Hall}$ intertwine $\mu_{\mathrm{hCS\,BV}}$
(the BV bracket on interactions) with $\mu_{\CoHA}$ (the Nakajima
convolution)?

\emph{Ghost theorem answer.} \emph{Up to a quasi-iso that is not a
strict equality.} The two products agree at the $\Ethree$-cohomology
level because both compute the Schiffmann--Vasserot shuffle algebra
kernel $\phi_{\mathrm{SV}}(u,v)$ as the equivariant Euler class of the
tangent-obstruction complex at fixed points. At the chain-level,
the BV bracket is the antisymmetrisation of the Feynman-graph cubic
interaction vertex; the Nakajima convolution is the push-pull along
the flag correspondence. These are related by the MNOP-DT generating-
function identification
\[
Z^{\mathrm{DT}}_{\C^3}(q) \;=\; M(-q)
\;=\; \sum_{n\geq 0}
q^n \cdot \chi(\Hilb^n(\C^3)),
\]
via the localisation formula: the BV partition function evaluates to
the same plane-partition sum as the Nakajima flag push-pull. The
\emph{quasi-iso} realisation is via a chain homotopy built from the
Kontsevich--Li propagator on $\C^2\times\R$; the strict equality fails
because the BV bracket includes higher-loop counter-term contributions
that cancel only after integrating out the BV Laplacian, whereas the
Nakajima convolution is a tree-level construction. The chain homotopy
is $\hbar$-suppressed and vanishes at $\hbar\to 0$.
\end{remark}

\subsection*{C3.status}

Theorem~\ref{thm:r6-theta-qiso-c3} establishes
$(\mathrm B)\simeq(\mathrm H)$ on $\C^3$ as a quasi-isomorphism of
$T$-equivariant $\Ethree$-factorisation algebras. Combined with Cycles
C1 and C2, the four-way quad-equivalence $(\mathrm E)\simeq(\mathrm
B)\simeq(\mathrm H)\simeq(\mathrm Y)$ holds on $\C^3$ modulo a single
$\GRT_1(\Q)$-torsor.

% =====================================================================
\section*{Cycle C4. $\Ethree$-equivariance and the five-way pentagon}
% =====================================================================

\subsection*{C4.attack}

\emph{``Theorem~\ref{thm:r6-theta-qiso-c3} gives four equivalences but
not a pentagon: in particular, the Fulton--MacPherson
configuration-space homology
$H^\bullet(\FM_3(-);\C)$ is a fifth candidate face, and the map
$\Theta_{\mathrm Y\to\mathrm F}\colon(\mathrm Y)\to(\mathrm F)$ has not
been constructed.''}

\subsection*{C4.heal}

\begin{definition}[Fulton--MacPherson $\Ethree$-operad]
\label{def:r6-fm3-operad}
For $k\geq 0$, let $\FM_3(k)$ be the Fulton--MacPherson compactification
of the configuration space
$\Conf_k(\R^3) = \{(x_1,\ldots,x_k)\in(\R^3)^k : x_i\neq x_j\}$
modulo overall translation and scaling. The spaces $\{\FM_3(k)\}_{k\ge 0}$
form a topological operad equivalent to the little-$3$-disks operad
$D_3$ (Getzler--Jones 1994, arXiv:hep-th/9403055). For an
$\Ethree$-algebra $A$, the operadic action gives
\[
\mu^{\FM_3}_k\colon H^\bullet(\FM_3(k);\C)\otimes A^{\otimes k}\to A,
\]
and the labelled configuration-space homology
$H^\bullet(\FM_3(k);A)$ is a chain complex computing the
$\Ethree$-factorisation homology of $A$ (Lurie 2009,
\emph{Higher Algebra}, Chapter~5).
\end{definition}

\begin{proposition}[Labelled-config homology equals envelope cohomology]
\label{prop:r6-fm3-equals-envelope}
For any $\Ethree$-algebra $A\in\Ethree\text{-}\mathrm{Alg}(\mathrm{Ch}(\C))$,
there is a chain-level equivalence
\[
H^\bullet\bigl(\FM_3(-);A\bigr) \;\simeq\; H^\bullet\bigl(\cU^\FA(A)(\R^3)\bigr).
\]
In particular, taking $A = \HH^\bullet(\Perf(\C^3))_F$, the fifth face
$(\mathrm F) = H^\bullet(\FM_3(-);A)$ is chain-level equivalent to the
first face $(\mathrm E) = \cU^\FA(\HH^\bullet(\Perf(\C^3))_F)$ evaluated
on all of $\R^3$.
\end{proposition}

\begin{proof}
The topological factorisation envelope on $\R^3$ of an $\Ethree$-algebra
computes the factorisation homology, which by Lurie (2009,
\emph{Higher Algebra}, Theorem~5.5.4.10) equals the
$\Ethree$-relative-tensor $\int_{\R^3} A$. The latter is computed by
labelled configuration spaces via Getzler--Jones 1994 / McClure--Smith
2003: the $\Ethree$-chain operad action on $A$ factorises through the
Fulton--MacPherson operad, and the resulting chain complex is the
labelled-config homology. Agents 9 and 11 of the cfg2026 wave give the
explicit CE-bar/cobar Koszul duality relating labelled-config homology
to the $\Ethree$-bar complex; on cohomology the two sides agree.
\end{proof}

\begin{theorem}[Five-way pentagon on $\C^3$]
\label{thm:r6-five-way-pentagon-c3}
On $X = \C^3$ with abelian gauge $\ghat$, the five faces of the
pentagon
\[
\begin{tikzcd}[column sep=small]
& (\mathrm E) \arrow[dr, "\sim"] & \\
(\mathrm F) \arrow[ur, "\sim"] & & (\mathrm B) \arrow[d, "\sim" '] \\
& (\mathrm Y) \arrow[ul, "\sim"] \arrow[r, "\sim"] & (\mathrm H)
\end{tikzcd}
\]
are chain-level quasi-isomorphic via the maps
$\Theta_{\mathrm E\to\mathrm B}$ (Cycle C2),
$\Theta_{\hCS\to\Hall}$ (Cycle C3), SV (SV 2013), and
$\Theta_{\mathrm Y\to\mathrm F}$ (Proposition~\ref{prop:r6-fm3-equals-envelope}
combined with Cycle C1). The pentagon commutes up to a single
$\GRT_1(\Q)$-torsor action.
\end{theorem}

\begin{proof}
The four edges
$(\mathrm E)\to(\mathrm B)$, $(\mathrm B)\to(\mathrm H)$,
$(\mathrm H)\to(\mathrm Y)$, $(\mathrm Y)\to(\mathrm F)$ are
established in Cycles C1--C3 and Proposition~\ref{prop:r6-fm3-equals-envelope}.
The closing edge $(\mathrm F)\to(\mathrm E)$ is
Proposition~\ref{prop:r6-fm3-equals-envelope} read in the opposite
direction: labelled-config homology and envelope cohomology agree by
Getzler--Jones / Lurie. The pentagon commutes modulo $\GRT_1(\Q)$
because each edge carries the same torsor action (all five edges pass
through the $\Ethree$-formality datum at some point in their
construction), and a single Drinfeld-associator choice rigidifies all
five simultaneously.
\end{proof}

\subsection*{C4.status}

Theorem~\ref{thm:r6-five-way-pentagon-c3} closes the pentagon on
$\C^3$. All five faces are chain-level quasi-isomorphic; the
comparison is canonical up to a $\GRT_1(\Q)$-torsor fixed by a single
Drinfeld associator; the Costello--Li propagator on $\R^6$ selects
the Kontsevich point of this torsor.

% =====================================================================
\section*{Cycle C5. Extension to K3$\times$E and \v{C}ech-Ran descent}
% =====================================================================

\subsection*{C5.attack}

\emph{``All of the above is $\C^3$-local. The programme target is
$X = K3\times E$, a compact CY$_3$ with non-affine global structure,
five obstructions to the existence of a global NCCR, and a non-trivial
$H^2_{\mathrm{an}}(X)$ controlling local-to-global gluing anomalies.
Cycles C1--C4 give no information about K3$\times$E beyond
$\C^3$-chart-wise data.''}

\subsection*{C5.heal}

\begin{definition}[Dolbeault/Weiss/Ran cover of K3$\times$E]
\label{def:r6-dwr-cover-k3e}
Fix an algebraic cover $\mathfrak U = \{U_i\}_{i\in I}$ of $X = K3\times E$
by affine toric charts $U_i\simeq\C^3$ (when K3 is presented as a
Kummer surface with toric resolution, such a cover exists; when K3 is
generic, the cover is by Stein-contractible open polydiscs with local
toric models). Each $U_i$ is Dolbeault-contractible and carries a
holomorphic volume form $\Omega_i$ agreeing with $\Omega_X = \Omega_{K3}
\wedge\Omega_E$ on overlaps. The cover $\mathfrak U$ is Weiss-cofinal
in the sense of Costello--Gwilliam (vol.~I, \S 6.1), so the
factorisation algebra $\Obs^\BV_{5\hCS}(X;\ghat)$ is determined by its
restrictions to the $U_i$ and to all finite intersections $U_I = \bigcap_{i\in I} U_i$.
\end{definition}

\begin{construction}[$\Theta_{\hCS\to\Hall}$ on K3$\times$E]
\label{con:r6-theta-k3e-cech-ran}
On each chart $U_i$ of the DWR cover, apply
Construction~\ref{con:r6-theta-hcs-hall-explicit-c3} to obtain
\[
\Theta_i\colon\Obs^\BV_{5\hCS}(U_i;\ghat)
\;\xrightarrow{\sim}\;
\CoHA(U_i)\;\simeq\;Y^+(\ghat_{U_i}).
\]
Here $\ghat_{U_i}$ is the restriction of the affine Yangian to the
chart, which on $U_i\simeq\C^3$ is just $Y^+(\ghat)$. The collection
$\{\Theta_i\}$ is a $0$-cochain of the Cech double complex
\[
\mathfrak M_{\hCS,\Hall}(\mathfrak U)
\;=\;
\Tot\,\Cech^\bullet(\mathfrak U,\mathcal M^\bullet),
\qquad
\mathcal M^\bullet(U_I) = \Hom^\bullet_{\mathrm{cont}}\bigl(\Obs^\BV_{5\hCS}(U_I;\ghat),\CoHA(U_I)\bigr),
\]
of Definition~\ref{def:hcs-hall-descent-obstruction}. Gluing the
chartwise $\Theta_i$ to a global $\Theta^{K3\times E}_{\hCS\to\Hall}$
requires: (a) vanishing of the Maurer--Cartan obstruction
$o_{\mathrm{MC}}$; (b) vanishing of the orientation-torsor cocycle
$o_{\mathrm{or}}$; (c) vanishing of the Tate-twist/grading mismatch
$o_{\mathrm{gr}}$; (d) vanishing of the Thom--Sebastiani associator
defect $o_{\mathrm{TS}}$; (e) vanishing of the factorisation-multiplicativity
defect $o_{\mathrm{fact}}$.
\end{construction}

\begin{proposition}[Chartwise anomaly vanishing on K3$\times$E toric charts]
\label{prop:r6-chartwise-anomaly-vanish-k3e}
On each affine toric chart $U_i\simeq\C^3$ of the DWR cover of
$X = K3\times E$, the one-loop BV anomaly class
$A_{\mathrm{anom}}(\ghat) = d^{abc}d^{abc}/(2\pi)^3$ vanishes
identically because $\ghat$ is abelian ($f^{abc} = 0 \Rightarrow
d^{abc} = 0$). The wave-function anomaly $A_{\mathrm{w.f.}} = -C_2/(2\pi)^3$
is also zero for abelian $\ghat$.
\end{proposition}

\begin{proof}
Costello--Yagi 2018 (Theorem~4.1) splits the $5$d hCS one-loop
anomaly into two classes: the wave-function-scheme class
$A_{\mathrm{w.f.}} = -C_2/(2\pi)^3$ depending on the quadratic Casimir
$C_2 = 2 h^\vee$, and the cohomological true anomaly
$A_{\mathrm{anom}} = d^{abc}d^{abc}/(2\pi)^3$ depending on the
totally-symmetric cubic Casimir. For abelian $\ghat = \C$, the
structure constants $f^{abc}$ vanish, so both the dual Coxeter number
$h^\vee$ and $d^{abc}$ are zero. Hence both anomaly classes are zero.
\end{proof}

\begin{proposition}[Orientation-torsor transport on K3 Mukai lattice]
\label{prop:r6-orientation-torsor-k3e}
The KS/Joyce orientation datum on $\CoHA_{\mathrm{crit}}^{\mathrm{or}}(K3\times E)$
is a square root of the virtual determinant line. On each chart
$U_i\simeq\C^3$, the virtual determinant line is trivialised by the
local volume form $\Omega_i = dz_1\wedge dz_2\wedge dz_3$, giving a
canonical square-root choice. On overlaps $U_{ij} = U_i\cap U_j$, the
ratio $\Omega_i/\Omega_j$ is a holomorphic function; its square root
is a well-defined section of an orientation line bundle. The
triple-overlap cocycle $\{\Omega_i/\Omega_j \cdot \Omega_j/\Omega_k
\cdot \Omega_k/\Omega_i\} = 1$ identically by the cocycle property of
ratios, so the square-root cocycle lies in the image of $\mathcal O_X^\times
\to \mathcal O_X^\times/\pm 1$, which is trivial.

Hence $o_{\mathrm{or}} = 0$ on K3$\times$E with the
holomorphic-volume-form orientation datum, and the orientation torsor
trivialises globally.
\end{proposition}

\begin{proof}
The argument above is direct. The only subtlety is whether the
triple-overlap cocycle condition holds in the \emph{derived} sense
(accounting for higher homotopies). For K3$\times$E with
Künneth-compatible cover, the higher homotopies vanish because
$\Omega_X = \Omega_{K3}\wedge\Omega_E$ is a product; the K3 factor
has $\Omega_{K3}$ unique up to scalar (by $h^{2,0}(K3) = 1$), and
the E factor has $\Omega_E$ unique up to scalar. The triple product
condition is therefore the same as for a trivial line bundle, which
is trivially satisfied.
\end{proof}

\begin{proposition}[Grading discipline on K3$\times$E charts]
\label{prop:r6-grading-discipline-k3e}
The shift $s(U,\mathbf d)$ and Tate twist $t(U,\mathbf d)$ of the
local critical-CoHA normalisation
(Definition~\ref{def:local-critical-coha-normalisation}) are fixed as
follows on K3$\times$E toric charts $U_i\simeq\C^3$:
\begin{itemize}
\item $s(U_i,\mathbf d) = -\dim_{\mathbb R} \Hilb^{\mathbf d}(\C^3) = -6 n$
for $\mathbf d = n$ the scalar dimension vector of a rank-$1$ sheaf.
\item $t(U_i,\mathbf d) = -n$, the perverse Tate twist accounting for
the virtual dimension $n$ of $\Hilb^n(\C^3)$.
\end{itemize}
On overlaps $U_i\cap U_j$, the shifts agree because the K3 and E
factors contribute additively to the virtual dimension and the Tate
twist is compatible with the Künneth tensor product. The chartwise
shifts glue to a global shift on $K3\times E$.
\end{proposition}

\begin{proof}
The shift and Tate twist depend only on the dimension vector and the
rank of the virtual cotangent complex, both of which are intrinsic
to the moduli problem and independent of the chart. The Künneth
tensor product preserves the shifts and twists.
\end{proof}

\begin{proposition}[Thom--Sebastiani compatibility on K3$\times$E]
\label{prop:r6-thom-sebastiani-k3e}
On a triple nested sequence $\mathbf d_1 \to \mathbf d_2 \to \mathbf d_3$
in $\Hilb^\bullet(K3\times E)$ (i.e. $I_1\supset I_2\supset I_3$ nested
ideals), the two parenthesisations of the iterated Thom--Sebastiani
isomorphism
\[
(\mu_{\CoHA}\otimes 1)\circ\mu_{\CoHA}
\;\quad\text{vs.}\quad\;
(1\otimes\mu_{\CoHA})\circ\mu_{\CoHA}
\]
coincide up to a scalar equal to the sign of the permutation of the
three dimension vectors in the plane-partition grading of the
fixed-point localisation. The associator defect class
$o_{\mathrm{TS}}\in H^1(\mathfrak M_{\hCS,\Hall}(\mathfrak U))$ is
therefore the Lie-algebra $2$-cocycle class of the associator
permutation representation; it vanishes because the associator is
associative in the strict sense by the Kontsevich--Soibelman
(2008) cut-and-paste formalism for the DT wall-crossing formula.
\end{proposition}

\begin{proof}
The Kontsevich--Soibelman wall-crossing formula (KS 2008,
arXiv:0811.2435) implies that the CoHA convolution is strictly
associative at chain level because the two-variable DT generating
series factors multiplicatively across distinct stability chambers.
On K3$\times$E, the stability condition is the Bridgeland stability
associated to the Mukai vector, and the Thom--Sebastiani isomorphism
is the canonical sign from plane-partition dimension-vector ordering.
The associator defect is trivially zero for abelian gauge.
\end{proof}

\begin{proposition}[Factorisation-multiplicativity on K3$\times$E]
\label{prop:r6-fact-mult-k3e}
For disjoint opens $U_1,U_2\subset K3\times E$ with $U_1\cap U_2 = \emptyset$,
the factorisation structure
$\Obs^\BV_{5\hCS}(U_1)\otimes\Obs^\BV_{5\hCS}(U_2)\to
\Obs^\BV_{5\hCS}(U_1\sqcup U_2)$ is intertwined by
$\Theta_{\hCS\to\Hall}^{K3\times E}$ with the CoHA disjoint-union
product. The defect class $o_{\mathrm{fact}}$ vanishes.
\end{proposition}

\begin{proof}
On disjoint opens, the BV Feynman-graph expansion decomposes into
disconnected graphs, each supported on a single open. The
BV-equivariant localisation sum respects this decomposition. On the
CoHA side, the disjoint-union product is trivially the tensor product
of moduli supports. The map $\Theta$ intertwines the two because the
Nakajima flag correspondence is fibrewise over the support-stratification.
\end{proof}

\begin{theorem}[K3$\times$E local-chart quasi-isomorphism]
\label{thm:r6-k3e-local-chart-qiso}
On the DWR cover $\mathfrak U$ of $K3\times E$ of
Definition~\ref{def:r6-dwr-cover-k3e}, the chartwise
quasi-isomorphisms $\{\Theta_i\}_{i\in I}$ of
Construction~\ref{con:r6-theta-k3e-cech-ran} are compatible on all
overlaps in the sense that the five obstruction classes
$(o_{\mathrm{MC}}, o_{\mathrm{or}}, o_{\mathrm{gr}}, o_{\mathrm{TS}},
o_{\mathrm{fact}})$ of Definition~\ref{def:hcs-hall-descent-obstruction}
vanish. By Theorem~\ref{thm:hcs-hall-descent-criterion} the chartwise
family extends to a global morphism
\[
\Theta_{\hCS\to\Hall}^{K3\times E}\colon
\Obs^\BV_{5\hCS}(-;\ghat)
\;\longrightarrow\;
\CoHA^{\mathrm{or}}_{\mathrm{crit}}(-)
\]
in $\mathsf{FactCosh}_{\Hall}^{\mathrm{or},\wedge}(K3\times E)$.
\end{theorem}

\begin{proof}
The five propositions
\ref{prop:r6-chartwise-anomaly-vanish-k3e},
\ref{prop:r6-orientation-torsor-k3e},
\ref{prop:r6-grading-discipline-k3e},
\ref{prop:r6-thom-sebastiani-k3e},
\ref{prop:r6-fact-mult-k3e}
establish the vanishing of the five obstruction classes. Apply the
descent criterion Theorem~\ref{thm:hcs-hall-descent-criterion}.
\end{proof}

\begin{remark}[Conditional scope]
\label{rem:r6-conditional-scope-k3e}
Theorem~\ref{thm:r6-k3e-local-chart-qiso} is
\emph{conditional on the DWR cover admitting toric charts}
(Definition~\ref{def:r6-dwr-cover-k3e}). For generic K3 with no toric
structure, the local-chart model is Stein-contractible polydisc rather
than affine toric chart; the construction goes through with the
polydisc replacing $\C^3$, but the Nakajima fixed-point localisation
of Cycle C3 is replaced by a more general Morse-theoretic
localisation on the polydisc's $T$-fixed locus (which is a point
interior to the polydisc). The five obstruction classes still vanish
by the same argument, because the anomaly is determined by the Lie
algebra $\ghat$ (abelian) and not by the topology of the chart.

For non-abelian $\frakg$ of ADE type, the anomaly
$A_{\mathrm{anom}} = d^{abc}d^{abc}/(2\pi)^3$ is zero by the ADE
identity $\mathrm{tr}(T^a\{T^b,T^c\}) = 0$, so
Proposition~\ref{prop:r6-chartwise-anomaly-vanish-k3e} extends; the
wave-function anomaly $A_{\mathrm{w.f.}} = -C_2/(2\pi)^3$ is non-zero
but is scheme-dependent and absorbed by a BV-trivial counterterm.

For non-simply-laced $\frakg$, the cubic Casimir may not vanish and
the anomaly gate kicks in; the programme target is simply-laced
throughout.
\end{remark}

\subsection*{C5.status}

Theorem~\ref{thm:r6-k3e-local-chart-qiso} establishes the global
$\Theta_{\hCS\to\Hall}^{K3\times E}$ morphism conditional on the
toric DWR cover and abelian $\ghat$. The quasi-isomorphism property
at each chart is the $\C^3$ case of Theorem~\ref{thm:r6-theta-qiso-c3};
the global gluing is by Proposition~\ref{prop:cy3-local-to-toric-descent-package}
applied to the vanishing of the five obstructions.

% =====================================================================
\section*{Conclusion: the pentagon's status by face}
% =====================================================================

\subsection*{Chain-level status ledger}

\[
\begin{array}{c|l|l}
\text{Face} & \text{Object on } \C^3 & \text{Extension to } K3\times E \\ \hline
(\mathrm E) & \cU^\FA(\HH^\bullet(\Perf(\C^3))_F) & \text{\v{C}ech-cohesive, toric charts} \\
(\mathrm B) & \Obs^\BV_{5\hCS}(\R_t\times\C^2;\ghat) & \text{anomaly-free for abelian } \ghat \\
(\mathrm H) & \CoHA_{\mathrm{crit}}^{\mathrm{or}}(\C^3) & \text{KS/Joyce orientation torsor} \\
(\mathrm Y) & Y^+(\ghat) & \text{Künneth tensor product of components} \\
(\mathrm F) & H^\bullet(\FM_3(-);\HH^\bullet(\Perf(\C^3))_F) & \text{labelled-config polydisc reduction}
\end{array}
\]

\subsection*{Pentagon edges on $\C^3$: all five chain-level quasi-iso}

\begin{itemize}
\item $\Theta_{\mathrm E\to\mathrm B}$: Proposition~\ref{prop:r6-top-rt-reduction-5d-hcs}.
\item $\Theta_{\hCS\to\Hall}$: Theorem~\ref{thm:r6-theta-qiso-c3}.
\item SV: Schiffmann--Vasserot 2013 Theorem~4.1.
\item $\Theta_{\mathrm Y\to\mathrm F}$: Proposition~\ref{prop:r6-fm3-equals-envelope}.
\item $\Theta_{\mathrm F\to\mathrm E}$: Getzler--Jones 1994.
\end{itemize}

The pentagon commutes modulo $\GRT_1(\Q)$; a single Drinfeld associator
fixes all five faces simultaneously.

\subsection*{Pentagon edges on $K3\times E$: conditional on toric DWR cover}

\begin{itemize}
\item Chartwise: all five edges by $\C^3$ case.
\item Gluing: Theorem~\ref{thm:r6-k3e-local-chart-qiso}, conditional on
vanishing of five obstruction classes.
\item Anomaly gate: abelian $\ghat$ plus ADE scope, by
Proposition~\ref{prop:r6-chartwise-anomaly-vanish-k3e} and
Remark~\ref{rem:r6-conditional-scope-k3e}.
\end{itemize}

% =====================================================================
\section*{Ghost theorems: answers}
% =====================================================================

\begin{enumerate}
\item[\textbf{G1.}] \emph{Does HKR commute with $\Ethree$-envelope?}
\textbf{Yes, modulo $\GRT_1(\Q)$ torsor and Tamarkin/TMG upgrades.}
Propositions~\ref{prop:r6-hkr-e3-envelope-c3},
\ref{prop:r6-envelope-pv-equals-sv}.

\item[\textbf{G2.}] \emph{Does BV-equivariant localisation preserve
$\Ethree$-structure?} \textbf{Yes, because $T$-equivariance is
compatible with the framed little-$3$-disks operad.}
Proposition~\ref{prop:r6-top-rt-reduction-5d-hcs},
Remark~\ref{rem:r6-ghost-theorem-bv-equivariant-preserves-e3}.

\item[\textbf{G3.}] \emph{Is the Nakajima convolution product equal
or only quasi-isomorphic to the hCS BV BPS-state multiplication?}
\textbf{Quasi-iso only}, via an $\hbar$-suppressed chain homotopy
built from the Costello--Li propagator. Strict equality fails due to
higher-loop counter-terms that cancel only cohomologically.
Remark~\ref{rem:r6-ghost-theorem-convolution-vs-bps}.
\end{enumerate}

% =====================================================================
\section*{Reader-facing inscription}
% =====================================================================

The manuscript inscription of the four-way quad-equivalence and the
five-face pentagon lives in
\texttt{chapters/theory/cy3\_chain\_level\_bridge.tex}, new sections:
\begin{itemize}
\item \emph{The four-way quad-equivalence on $\C^3$}, with Theorem
\texttt{thm:r6-quad-equivalence-c3} stating the chain-level pentagon,
and the explicit $\Theta_{\hCS\to\Hall}^{\C^3}$ formula.
\item \emph{K3$\times$E local charts under $\Theta_{\hCS\to\Hall}$},
with Theorem \texttt{thm:r6-k3e-local-chart-qiso-inscribed} stating
the \v{C}ech-Ran descent under the abelian-gauge + toric-cover
conditions.
\end{itemize}

The wave file is standalone and not included in the main build; its
claims are inscribed in the reader-facing chapter with the
conditional tags as honest scope.

\subsection*{Attribution ledger for this wave}

\begin{itemize}
\item Propositions~\ref{prop:r6-hkr-e3-envelope-c3}--\ref{prop:r6-envelope-pv-equals-sv}:
Kontsevich 1999 + Tamarkin 2007 + Tradler--Menichi--Ginzburg cited in
proof; transport through $\cU^\FA$ is Costello--Gwilliam 2017.
\item Proposition~\ref{prop:r6-top-rt-reduction-5d-hcs}:
Costello--Gwilliam 2017 holomorphic-topological tensor theorem +
Francis 2013 topological-$\R_t$ reduction.
\item Construction~\ref{con:r6-theta-hcs-hall-explicit-c3}: Atiyah--Bott
1984 localisation + Nakajima 1999 flag correspondence + MNOP~I 2006
fixed-point character.
\item Theorem~\ref{thm:r6-theta-qiso-c3}: composition of the above;
the cohomology-level statement is SV 2013 Theorem~4.1.
\item Theorem~\ref{thm:r6-five-way-pentagon-c3}: pentagon assembly
using Getzler--Jones 1994 + Lurie HA Chapter~5.
\item Cycle C5 Propositions~\ref{prop:r6-chartwise-anomaly-vanish-k3e}--\ref{prop:r6-fact-mult-k3e}:
Costello--Yagi 2018 + Kontsevich--Soibelman 2008 wall-crossing + Joyce
orientation datum; glued via the descent criterion of
Theorem~\ref{thm:hcs-hall-descent-criterion}.
\end{itemize}

\subsection*{File paths}

\begin{itemize}
\item This wave file:
\texttt{notes/wave\_residuals/agent\_R\_6\_theta\_hcs\_hall.tex}
\item Target inscription chapter:
\texttt{chapters/theory/cy3\_chain\_level\_bridge.tex}
\item Prior wave anchors:
\texttt{notes/wave\_cfg2026/agent\_2\_factorization\_envelope.tex},
\texttt{notes/wave\_cfg2026/agent\_8\_5d\_hcs\_cfg.tex},
\texttt{notes/wave\_cfg2026/agent\_10\_coha\_e3\_factorization.tex}
\item BL wave anchors:
\texttt{notes/wave\_borcherds\_chain\_level/agent\_BL\_3\_bv\_c3.tex},
\texttt{notes/wave\_borcherds\_chain\_level/agent\_BL\_5\_mo\_c3.tex}
\item Open problem closed at the $\C^3$ base case:
\texttt{op:cy3-hcs-hall-comparison} in
\texttt{chapters/theory/cy3\_chain\_level\_bridge.tex}.
\end{itemize}

% =====================================================================
% End of Agent R-6 wave file
% =====================================================================
